\section{Related Works}
\label{sec:related_works}

\noindent
Active Learning and Radiance Fields are prosperous research fields with various research directions. In this section, we limit our literature review to works that directly relate to our tasks. We refer the readers to literature reviews if they are interested in Radiance Fields~\cite{gao2023nerf-review, dellaert2021neural-review} or Deep Active Learning~\cite{ren2021survey,zhan2022comparative}.
Besides, Fisher Information has been extensively studied in deep active learning~\cite{Ash2021GoneFishing,Ash2020Deep,kothawade2021similar,kirsch2021stochastic}.
Notably, Kirsch~\etal~\cite{kirsch2022unifying} unified existing works in active learning for deep learning problems from the perspective of Fisher Information, which shares many insights with our method. 
\paragraph{Uncertainty Quantifications for Radiance Fields}
In the Neural Radiance Fields, Lee~\etal~\cite{lee2022uncertainty}, Yan~\etal~\cite{yan2023activeIO} and Zhan~\etal~\cite{zhan2022activermap} attempted to quantify the uncertainty in a scene by the distribution of densities on a casted ray. Shen~\etal~\cite{CF-NeRF,shen2021snerf} designed Bayesian models by re-parametrizing the NeRF model. However, they only tackled the predictive uncertainty of the model that did not relate to the observed information of the parameters. 
Sunderhauf~\etal\cite{sünderhauf2022densityaware} propose an additional uncertainty measure in uncharted regions, determined by the ray termination probability on the learned geometry.
Regarding concurrent work, Goli~\etal~\cite{goli2023} introduces a hypothetical perturbation field and interpolates uncertainty over NeRF-style query points for uncertainty quantification and NeRF clean-up task.
The perturbation field with respect to query points is also not compatible with recent radiance field models like 3D Gaussian Splatting~\cite{kerbl3Dgaussians}, which do not take query points as inputs.
Besides, their algorithm is not designed for an active learning problem since selecting candidates with the largest per-pixel uncertainty may not be optimal, as the information gain for the model parameters is not concerned.
In contrast, we focus on active learning problems. We derive our theory from the expected information gain for the model and compute the Hessian for the model parameters, allowing us to apply our algorithm to recent advancements in radiance fields such as 3D Gaussian splatting. 
Our paper focuses on active view selection and mapping based on our novel expected information gain objective. 
Moreover, their implementation relies on multiple function calls to PyTorch backward engine whereas our efficient CUDA implementation for the diagonal Hessian computation consumes much less time (11.3 ms $\pm$ 33.3 $\mu$s v.s. 1.1 s $\pm$ 2.5 ms).
To the best of our knowledge, we are the first method that quantifies the uncertainty directly on model parameters for Radience Fields, thanks to the powerful framework of 3D Gaussian and our efficient formulation. 
\paragraph{Active View Selection and Mapping}
Although the next best view selection problem was extensively studied before radiance fields gained popularity, there has not been much literature on active ``training'' view selection for volumetric rendering. 
Jin~\etal~\cite{jin2023neu} trains a neural network to predict per-pixel uncertainty for view selection. However, their method requires an image capture as the input of the network, which is closer to dataset subsampling rather than active learning.
ActiveNeRF~\cite{pan2022activenerf} studied the next best view selection by directly adding variances as the output of the vanilla NeRF, and it is the closest approach to the view selection problem we are working on. 
In a larger scope, Active Perception~\cite{bajcsy1988active,bajcsy2018revisiting} has been studied along with the advancement of robotics and computer vision.
Traditionally, frontier-based approach~\cite{yamauchi1997frontier} and Rapidly exploring Random Tree~(RRT)~\cite{lavalle1998rrt} are the most representative and have been extended and enhanced for different data representations~\cite{zhu20153d, karaman2010incremental, dornhege2013frontier, shen2012autonomous} and complex conditions~\cite{ye2022multi, papachristos2017uncertainty}.
For the larger scope of Active Simultaneous Localization and Mapping~(SLAM), please refer to literature reviews~\cite{lluvia2021active-survey, placed2023survey}.
In recent years, Guedon~\etal~\cite{guedon2022scone,guedon2023macarons} attempted to improve the coverage for mapping systems using point cloud representation with neural networks. Dhami~\etal~\cite{dhami2023pred} studied the next best view selection in view planning by first completing the point cloud from partial observation. Georgakis~\etal~\cite{georgakis2022uncertainty} measured uncertainty through deep ensemble. Chaplot~\etal~\cite{chaplot2020learning} and Ramakrishnan~\etal~\cite{ramakrishnan2020occupancy} developed active exploration methods with deep reinforcement learning.
The mapping system that uses implicit feild~\cite{zhu2021niceslam, Sucar:etal:ICCV2021, Ortiz:etal:iSDF2022, sandstrom2023point} espeically 3D Gaussians~\cite{keetha2023splatam, yan2023gs, Matsuki:Murai:etal:CVPR2024,yugay2023gaussianslam} as representation has been widely studied.
This presents the need for active mapping algorithms that could maximize the mapping efficiency in an embodied robot system.
Similar to ActiveNeRF~\cite{pan2022activenerf}, Ran~\etal~\cite{Ran2023neurar} performed trajectory planning for object-centric autonomous reconstruction systems by predicting variance as an extra output of the network. 
Yan~\etal\cite{yan2023active-neural-mapping} discussed the intuition of flat minimum and quantified the predictive uncertainty of neural mapping models through the lens of loss landscape. However, they only approximate the uncertainty of the variances from the output of neighboring points. Our method directly quantifies the observed information of the parameters by computing the Hessian matrix of the log-likelihood function of the model's objective.
Besides, we introduce the Expected Information Gain for mapping systems from the first principle as a complementary method to previous heuristics. 
